\section{Conclusion and Discussion}
We have proposed a novel model of movement primitives based on the motion manifold hypothesis, named Motion Manifold Primitives++ (MMP++). This model synergizes the strengths of the existing MMP with those of conventional parametric curve representation-based primitive models, achieving enhanced motion generation accuracy and adaptability. Moreover, to overcome the geometric distortion issue in MMP++, we have introduced a CurveGeom Riemannian metric for the parametric curve space and presented Isometric Motion Manifold Primitives++ (IMMP++). This approach ensures that the latent space preserves the geometric structure of the motion manifold, which, in some instances, leads to significantly improved density fitting results within the latent space.

We highlight that the low dimensionality of the latent parametrization of the motion manifold has led to multiple advantages. This allows us to exploit non-parametric density fitting techniques, such as kernel density estimation, in the latent space, which are typically challenging to apply in high-dimensional spaces, for example, directly in the curve parameter space ${\cal W}$. Additionally, re-planning has been easily accomplished by finding a new latent value $z$, formulated as an optimization problem that is efficiently solvable thanks to the low dimensionality of the latent space.

We have extended MMP++ to accommodate matrix Lie group data, specifically SO(3), by designing a parametric curve model that utilizes the exponential map, logarithm map, and group action. Given that the water-pouring SE(3) trajectory dataset forms a connected manifold without a clustering structure, MMP++ -- even without isometric regularization -- has shown satisfying performance. However, in instances where we encounter a dataset with clustering structures, where a distorted latent space could result in poor density fitting, as verified in the other 2-DoF and 7-DoF robot experiments, employing isometric regularization alongside an appropriate Riemannian metric can offer a solution. Although developing isometric regularization for matrix Lie group data falls outside the scope of our paper, it represents an interesting direction for future research.

Although our focus has been on via-point models, we can adopt other parametric curve models with stronger inductive biases. Even more sophisticated movement primitives, such as those based on dynamical systems, can be used, for example, by making the parameters of the dynamical systems the outputs of the neural network, as done in~\cite{bahl2020neural}. These can potentially be combined with the autoencoder-based manifold learning technique. These combinations will yield a variety of manifold-based motion primitives capable of generating diverse motions, demonstrating high adaptability in previously unseen dynamic environments.

Lastly, we note that our current motion manifold primitives are not conditioned on vision inputs such as images, point clouds, or sequences of images, so the learned manifold of trajectories is applicable only to the trained environment, such as the specific shelf in Fig.~\ref{fig:manifol_results}. If the MMP model could generalize to environments it has not encountered during training, it would be extremely powerful. For example, if arbitrary shelf geometries were given as vision inputs, and the model could generate a manifold of diverse trajectories conditioned on those inputs, it would significantly enhance its applicability. Adopting techniques (e.g., network architectures) from recent LfD methods for visuomotor policy learning~\cite{chi2023diffusion, lee2024behavior} is a feasible direction to extend our framework to generate a manifold of trajectories conditioned on vision inputs. For example, the decoder can be modified to take a vision feature \(y\) along with the latent variable \(z\), so that the decoder maps \((z, y)\) to the curve parameter \(w\). Although this would require a relatively large amount of data paired with visual observations, our core idea remains applicable.
